\section{本章小结}
\label{sec:dtsp-summary}

本章开头提出了几个问题：服务器从正常状态出发，若干步后拥堵的概率是多少；赌徒能否靠下注或停止规则改变期望收益；随时查看的累计误差最坏有多大；相关的采样序列怎样给出可信的积分；带噪更新怎样收敛到所需的根。回答这些问题用到了三类结构，它们回答不同的问题，可以同时出现在同一个过程上。

第一类是转移结构。一步转移核决定全部有限维分布，多步分布由核的复合给出（定理~\ref{thm:dtsp-chapman-kolmogorov}）。状态能否到达、是否返回、返回是否够快，分别由沟通类、常返性与平均返回时间刻画；不变分布是分布更新的不动点，细致平衡是验证它的便利充分条件。存在不变分布、它唯一、从任意初态收敛到它，是三件不同的事：有限不可约链保证前两件，加上非周期才保证第三件，并且收敛是几何的（定理~\ref{thm:dtsp-finite-convergence}）。

第二类是条件均值结构。Doob分解把任何适应累计量拆成可预测漂移与鞅（定理~\ref{thm:dtsp-doob-decomposition}）。可预测下注不能改变公平游戏的期望；有界停时的可选抽样保持期望，无界停时则需要控制、有限平均时间或一致可积等条件，加倍下注说明这些条件缺一不可。固定预算下，Doob最大不等式、Azuma与Freedman界控制整段运行的最大偏差，后者按实际的条件方差计费；向量自归一化界把同样的方法推广到可预测设计。无穷时间上，$L^1$有界的鞅几乎必然收敛，一致可积再保证期望保持，$L^2$有界给出均方收敛；有限时间的概率界与无穷时间的收敛是两类不同的结论。

第三类是长期平均。平稳性描述时间平移下规律不变；遍历定理说明不可约正常返链的时间平均收敛到不变分布下的平均，它不需要非周期性。相关样本的平均误差由积分自相关时间放大，中心极限定理借助Poisson方程把相关和归结为鞅。同一个公平随机游走同时是Markov链、鞅，并具有平稳增量，而它本身并不平稳，这个例子说明三类结构并不互斥。

两个应用检验了这些工具。采样积分中，独立Monte Carlo的误差按$n^{-1/2}$下降，方差与生成成本决定实际计算量；重要性抽样、控制变量、条件平均、对偶与分层在保持积分目标的同时改变方差，各有前提与代价；Metropolis--Hastings、Gibbs与Hamilton蒙特卡洛用细致平衡或保体积对合保持目标分布，再由遍历定理把时间平均当作积分，误差评估必须计入相关性；路径重加权的似然比是均值为一的鞅，长路径时权重退化。带噪求根中，步长和发散与平方可和分别保证到达与抑制噪声，几乎超鞅定理给出线性与强单调情形的完整证明，平均ODE方法在有界性与漂移正则性下把随机轨迹与确定性稳定性连接起来，Poisson方程处理Markov噪声，快慢尺度与迭代平均分别处理耦合更新与波动。

使用这些结论时，应分别核验三件事：递推或采样链本身是否稳定（有界、回返、混合），数值计算是否收敛到所设定的目标（步长与可积条件、遍历与方差条件），以及这个目标是否就是统计上需要的量（模型与识别）。前两件是本章的内容，第三件属于统计推断。下一章把时间改为连续：转移核变成半群与生成元，停止时刻可以取不可数多个值，噪声的累计需要随机积分，这些问题保留本章的结构，但需要新增路径正则性与可积性条件。
